\chapter*{Before you start}
\addcontentsline{toc}{chapter}{Before you start}
\markboth{\textsc{Before you start}}{}

Most project tutorials open with ``first, set up your environment'' and then
spend one paragraph on something that will actually take you two evenings and
teach you a dozen things you will use for the rest of your life.

This book is that paragraph, expanded honestly.

\section*{What you will have at the end}

\begin{itemize}
  \item A Raspberry Pi you can reach by typing \code{ssh brewer} --- no
        password, no IP address to remember, from anywhere on your network.
  \item A Java project that compiles on your laptop and runs on the Pi one
        second later.
  \item Your code in git, backed up on GitHub, pushed with one command.
  \item An NFC reader wired up and answering.
\end{itemize}

None of that is specific to this project. Change the word ``NFC'' and you have
the setup for any embedded Linux work you will ever do.

\section*{Who this is for}

Someone who has written some C, can use a text editor, and has never done any of
the following: used SSH, generated a key pair, used git beyond \code{git clone},
run Maven, or set up a headless Linux machine.

If you have done all of those, read Chapter~\ref{ch:need} for the hardware
choices, skim Appendix~\ref{app:cheat}, and go straight to Volume One.

\section*{How this book is different from a wiki page}

\textbf{Everything is explained, not just listed.} You will not be told to run
\code{ssh-keygen -t ed25519} and left to wonder what happened. You will be told
what a key pair \emph{is}, why the private one never moves, and what would
actually go wrong if you skipped it.

That matters because setup steps are exactly the ones you cannot debug by
reading your own code. When SSH refuses you at midnight, understanding
\emph{why} it wants what it wants is the difference between five minutes and an
evening.

\begin{concept}{Checkpoints}
Each chapter ends with a \textbf{checkpoint}: a command to run, and the output
you should see.

\textbf{Do not continue past a failed checkpoint.} Setup problems compound
viciously --- a half-configured SSH key produces errors three chapters later
that look nothing like SSH.
\end{concept}

\section*{Conventions}

Commands are shown with a comment saying \emph{where} to run them, because half
of all setup confusion is running the right command on the wrong machine:

\begin{shell}
# on your laptop
ssh brewer

# on the Pi
ls -l /dev/spidev0.0
\end{shell}

\begin{itemize}
  \item \code{\$USER}, \code{martin}, \code{brewer.local} --- replace with your
        own username and hostname. The book uses \code{martin} and
        \code{brewer} throughout; pick your own in Chapter~\ref{ch:flashing}
        and stay consistent.
  \item \textcolor{brewerred}{Red rules} mark traps that have actually caught
        people. \textcolor{brewerblue}{Blue} explains a concept.
        \textcolor{breweramber}{Amber} is an analogy.
        \textcolor{brewergreen}{Green} is a checkpoint.
\end{itemize}

\section*{If you get stuck}

Appendix~\ref{app:trouble0} is indexed by \textbf{what you see on screen}, not by
topic. Search for your error message there first --- almost every setup failure
in this book has a known cause and a one-line fix.

\vspace{1em}
\begin{flushright}\itshape
Nothing here is hard. Some of it is just unfamiliar.
\end{flushright}
